% ============================================================
%  Preface
% ============================================================

\epigraph{%
  A theory is all the more impressive the simpler its premises,
  the more varied the things it relates,
  and the wider the scope of its applicability.%
}{Albert Einstein, \textit{Autobiographical Notes}}

This monograph began with a single ion.

A charged molecule --- small enough that quantum mechanics unambiguously applies,
large enough that Newtonian trajectory calculations reproduce its flight time
to within two parts per million --- was traced through four different mass
spectrometric platforms.
At each stage of its journey, we wrote down both descriptions.
The classical calculation and the quantum mechanical calculation gave the same
answer.

That agreement is the subject of this book.

It is not the trivial agreement of the \emph{correspondence principle}, which
asserts that quantum mechanics reduces to classical mechanics in some appropriate
limit of large quantum numbers.
It is a stronger, more precise statement: the two descriptions agree exactly,
at every mass measurement, at every energy, because both are consequences of
the same underlying mathematical structure.
That structure is \emph{bounded phase space}.

The Bounded Phase Space Law, stated formally in Part~II, is an empirical
observation: every persistent dynamical system occupies a bounded region of
phase space of finite Liouville measure.
From this single observation --- not a postulate, not a theoretical construct,
but a fact that every experiment since Galileo has confirmed --- we derive, in
sequence, the quantisation of energy, the four quantum numbers, the Pauli
exclusion principle, the Lorentz invariance of dynamics, the rest mass formula
$m_0 = \hbar\omega_0/c^2$, the equivalence of inertial and gravitational mass,
the emergence and quantisation of electric charge, the ideal gas law for any $N$,
the universal transport formula, the nuclear magic numbers, and
the acceleration of the cosmic expansion.

The monograph is organised around a metaphor: two crowns.

The \textbf{First Crown} is classical mechanics --- the crown of Newton,
Lagrange, and Hamilton.
Its language is trajectories, forces, potentials, and symmetries.
It has been the language of physics for three centuries, and it works.

The \textbf{Second Crown} is quantum mechanics --- the crown of Bohr,
Heisenberg, Schrödinger, and Dirac.
Its language is state vectors, operators, eigenvalues, and probability
amplitudes.
It was created to describe the phenomena that the First Crown could not, and
it too works --- often far better.

The two crowns sit on different heads, speak different languages, and
arrive at the same court.
The purpose of this monograph is to show that they do not merely agree
in the overlap of their domains; they are the same structure seen from
two vantage points.
The union is proved formally in Part~V, as the Triple Equivalence:
$\Osc \cong \Cat \cong \Part$, where $\Osc$ is the oscillatory (classical)
category, $\Cat$ is the categorical (information-theoretic) category, and
$\Part$ is the partition (quantum) category.

A third strand runs beneath both crowns: the \textbf{S-entropy calculus},
a scale of distance-to-truth defined on a bounded interval $[0, 100]$ relative
to a specific receiver.
S-entropy is not a heuristic; it is an axiomatic system with a Floor Theorem
(no bounded receiver attains perfect knowledge) and explicit functors connecting
it to both classical and quantum descriptions.
The mass spectrometric applications --- the generative spectral database, the
force-free ion trajectory, the GPU observation apparatus --- are all
instantiations of S-entropy on the instrument's bounded measurement space.

\bigskip
\noindent\textbf{How to read this book.}
Part~I (Chapters~1) is the entry point: a single ion's journey told twice,
in classical and quantum language simultaneously.
Readers who are already convinced of the framework's physical motivation may
begin at Part~II.
Readers who care primarily about the mass spectrometric applications may read
Part~I, skim Parts~II--IV for the key results, and then proceed directly to
Part~VIII.
The mathematical heart of the monograph is Part~V (Triple Equivalence);
the most comprehensive collection of derived results is Part~IX
(Empirical Confirmation).

\bigskip
\noindent\textbf{On style and notation.}
Physical derivations use the environment ``\textbf{Consequence}'' rather than
``Theorem''.
This choice is deliberate: the results are not discoveries but catalogue entries
--- the logical content of the Axiom made explicit.
A full list of macros is given in Appendix~B.

\bigskip
\noindent\textbf{Acknowledgements.}
The framework developed here emerged from years of work at the intersection of
cheminformatics, mass spectrometry, and theoretical physics.
[Acknowledgements to be completed.]

\vfill
\begin{flushright}
\textit{K.F.S.}\\
Munich, \today
\end{flushright}
